\section{Measuring coordination}\label{sec:measures}

\paragraph{Tracking and commitment.} Tracking is the mean share of attention the other agents put on the current true
source, averaged over the second half of a run: $1$ when everyone listens only to the source, $1/(N-1)$ when attention
is spread evenly. Commitment is the mean over agents of their largest single share, whoever it goes to. High commitment
with low tracking is the signature of confident betting on the wrong agent.

\paragraph{Benefit: attention as a bet.} Each agent's attention column is a bet on where the next news will appear.
When it appears at $s$, agent $i$ is charged the log score $-\ln A_{si}$; this log-loss is $\mathrm{KL}(\delta_s\,\|\,A_{\cdot
i})$ for a point mass on $s$. Uniform attention pays $\ln(N-1)$. The benefit is
\begin{equation}
  B=\ln(N-1)-\Big\langle -\ln A_{s_t^{\mathrm{true}} i}\Big\rangle_{i\ne s,\,t},
\end{equation}
at most $\ln(N-1)$ and strongly negative for confident bets on the wrong agent. It is the Kelly regret of the attention
bet and the quantity used by \citet{lee2024constructing}.

\paragraph{Information cost and calibration.} The information cost
$I=\langle\mathrm{KL}(A_{\cdot i}\,\|\,\mathrm{uniform})\rangle_i=\ln(N-1)-\langle H(A_{\cdot i})\rangle$ measures how much
an agent committed, not whether it was right: in the example of \cref{fig:cost-benefit} the right and the wrong bet both
cost $1.13$ nats. For a calibrated bettor the expected log-loss equals the entropy, so the calibration gap $B-I$ is near
zero; confident and wrong drives it far below zero.

\paragraph{Structural cost: attention as pipes.} Read the mutual attention $C_{ij}=(A_{ij}+A_{ji})/2$, scaled to sum to
one, as conductances. When news appears at $s$ it must reach everyone: inject one unit at $s$, withdraw $1/(N-1)$ at
every other agent, and let Kirchhoff's laws decide the currents $I_e$. The broadcast dissipation $\sum_e I_e^2/C_e$,
averaged over sources and divided by the same quantity for uniform attention, is the structural cost: $1$ for no
structure, about $2$ for a strong hub. A hub makes news from the hub cheap (spokes in parallel) but news from anyone else
expensive, since it must squeeze through one spoke. The measure is borrowed from adaptive transport networks, where it is
the quantity being minimized \citep{corson2010fluctuations}; the attention model never pays it, and we use it to price
what the model builds.

\figw{cost-benefit}{\linewidth}{Benefit and structural cost on a 7-agent example. Top: one agent's attention over the
other six ($s$ current source, $o$ old source) and what each bet earns. Bottom: the attention read as pipes (grey width
= conductance) carrying a broadcast from $s$ (coloured width = current). The hub holds $0.85$ of every other agent's
attention.}{cost-benefit}

\paragraph{Why not distance.} The original model judged networks by mean walk distance. Distance from $s$ depends on how
$s$'s audience is split, not on how much each listener gives to $s$: a network where everyone listens only to $s$ and one
where everyone listens a little to everyone have the same mean distance from $s$ (13.09 for every setting in E1). The
all-pairs version behaves erratically under noise (2.4, 6.0, 1.9 and 0.1 for $\nu=0$ to $0.75$ at $N=10$). Log-loss sees
concentration and is used throughout.

\paragraph{Shape: hierarchy coordinates.} Treeness $T_h\in[-1,1]$ (do paths converge on few top nodes?),
feedforwardness $F\in[0,1]$ (how much path weight avoids cycles?) and orderability $O\in[0,1]$ (how many nodes lie
outside cycles?) place a network in the hierarchy morphospace \citep{corominas2013originsheirarchy}. Weighted networks
are thresholded at 8 levels and the coordinates averaged, so a perfect star reads $7/8=0.88$. We compute treeness with
matrix powers in the path counts (eq.~14 of the original), correcting two errors in the code used by the 2020 thesis
(\cref{app:errata}). \Cref{fig:shapes} shows typical values.

\figw{shapes}{\linewidth}{Three shapes and their hierarchy coordinates (typical runs).}{shapes}
